\documentclass[10pt]{article}
\usepackage{geometry}
\geometry{
letterpaper,
margin={1in}
}

% required packages
\usepackage[hidelinks]{hyperref}
\usepackage{amsfonts}
\usepackage{amssymb}
\usepackage{amsmath}
\usepackage{amsthm}
\usepackage{mathtools}

% custom commands
\newcommand{\ket}[1]{\vert #1 \rangle}
\newcommand{\bra}[1]{\langle #1 \vert}
\newcommand{\braket}[2]{\langle #1 \vert #2 \rangle}
\DeclarePairedDelimiter\abs{\lvert}{\rvert}

% title
\title{CSCI 5244: Homework set 1}
\author{Rey Stone}
\date{}

\begin{document}

\maketitle

\begin{enumerate}
  \item Given a state of a qubit $\ket{\psi} = a \ket{0} + b \ket{1}$, where $a, b \in \mathbb{C}$ and $\abs{a}^2 + \abs{b}^2 = 1$, and recall that states $\ket{+} = \frac{1}{\sqrt{2}} (\ket{0} + \ket{1})$ and $\ket{-} = \frac{1}{\sqrt{2}} (\ket{0} - \ket{1})$

    \begin{enumerate}
      \item In 2-dimensional complex space with basis ${\ket{0}, \ket{1}}$, write down the matrix form of operator $A = \ket{+} \bra{+} + \frac{1}{2} \ket{-} \bra{-}$ and verify A is an observable (i.e., self-adjoint).

        First, the column vector of $\ket{0}$ can be written as $\begin{pmatrix} 1 \\ 0 \end{pmatrix}$ and $\ket{1}$ as $\begin{pmatrix} 0 \\ 1 \end{pmatrix}$.

        Solving for $\ket{+}$:

        \begin{align*}
          \frac{1}{\sqrt{2}} \left(\begin{pmatrix} 0 \\ 1 \end{pmatrix} + \begin{pmatrix} 1 \\ 0 \end{pmatrix} \right) = \frac{1}{\sqrt{2}} \begin{pmatrix} 1 \\ 1 \end{pmatrix}
        \end{align*}

        Solving for $\ket{-}$:

        \begin{equation*}
          \frac{1}{\sqrt{2}} \left(\begin{pmatrix} 0 \\ 1 \end{pmatrix} - \begin{pmatrix} 1 \\ 0 \end{pmatrix} \right) = \frac{1}{\sqrt{2}} \begin{pmatrix} 1 \\ -1 \end{pmatrix}
        \end{equation*}

        Corresponding $\bra{+}$ and $\bra{-}$ are just the complex conjugate of their respective kets in row vector form.
        
        The complex conjugate to both of these are represented as follows:

        \begin{align*}
          \bra{0} &= \frac{1}{\sqrt{2}} \begin{pmatrix} 1 & 1 \end{pmatrix} \\
          \bra{1} &= \frac{1}{\sqrt{2}} \begin{pmatrix} 1 & -1 \end{pmatrix} 
        \end{align*}

        Solving first operation $\ket{+}\bra{+}$:

        \begin{equation*}
          \frac{1}{\sqrt{2}} \begin{pmatrix} 1 \\ 1 \end{pmatrix} \cdot \frac{1}{\sqrt{2}} \begin{pmatrix} 1 & 1 \end{pmatrix} = \frac{1}{2} \begin{pmatrix} 1 & 1 \\ 1 & 1 \end{pmatrix}
        \end{equation*}

        Solving operation $\ket{-}\bra{-}$:

        \begin{equation*}
          \frac{1}{\sqrt{2}} \begin{pmatrix} 1 \\ -1 \end{pmatrix} \cdot \frac{1}{\sqrt{2}} \begin{pmatrix} 1 & -1 \end{pmatrix} = \frac{1}{2} \begin{pmatrix} 1 & -1 \\ -1 & 1 \end{pmatrix}
        \end{equation*}

        Matrix form of operator A:

        \begin{align*}
          A &= \frac{1}{2} \begin{pmatrix} 1 & 1 \\ 1 & 1 \end{pmatrix} + \frac{1}{2} \left( \frac{1}{2} \begin{pmatrix} 1 & -1 \\ -1 & 1 \end{pmatrix} \right) \\
            &= \frac{1}{2} \begin{pmatrix} 1 & 1 \\ 1 & 1 \end{pmatrix} + \frac{1}{4} \begin{pmatrix} 1 & -1 \\ -1 & 1 \end{pmatrix} \\
            &= \frac{1}{4} \begin{pmatrix}3 & 1 \\ 1 & 3 \end{pmatrix}
        \end{align*}

        Matrix $A$ is self-adjoint iff $A^{\dagger} = A$ where $A^{\dagger}$ is the conjugate transpose:

        \begin{equation*}
          A = \frac{1}{4} \begin{pmatrix} 3 & 1 \\ 1 & 3 \end{pmatrix}
        \end{equation*}

        Switch off-diagonal entries:

        \begin{equation*}
          \frac{1}{4} \begin{pmatrix} 3 & 1 \\ 1 & 3 \end{pmatrix}
        \end{equation*}

        Since we are not dealing with any imaginary numbers, no conjugation is needed. Thus proves that $A^{\dagger} = A$ making it self-adjoint, and making it a valid observable.

      \item If you measure the observable $A$ on the state $\ket{\psi}$, what are the possible measurement outcomes and associated probabilities?

        The measurement of observable $A$ "snaps" onto one of $A$'s eigenvalues $\ket{n}$ such that:

        \begin{equation*}
          Prob(a_n) = ||E_n \ket{\psi}||^2 = \bra{\psi} E_n \ket{\psi} \\
          \textit{\ (from the slides)}
        \end{equation*}

        The observable $A$ is already in the form $A = \ket{+}\bra{+} + \frac{1}{2} \ket{-}\bra{-}$, so the expected value of measurement can be written as:

        \begin{equation*}
          A = \sum_n \lambda_n \ket{v_n}\bra{v_n}
        \end{equation*}

        With possible outcomes being $\lambda_n$ with $Prob(\lambda_n) = |\braket{v_n}{\psi}|^2$.

        Already we can see that the two potential eigenvalues are $1$ in front of $\ket{+}\bra{+}$ and $\frac{1}{2}$ in front of $\ket{-}\bra{-}$

        The first probability Prob(1) can be calculated as such:

        \begin{align*}
          \bra{+} &= \frac{1}{\sqrt2}(\bra{0} + \bra{1}) = \frac{1}{\sqrt{2}} (( 1 \ 0 ) + ( 0 \ 1 )) = \frac{1}{\sqrt2} ( 1 \ 1 ) \\
         \braket{+}{\psi} &= \frac{1}{\sqrt2} (1 \ 1) \begin{pmatrix} a \\ b \end{pmatrix} = \frac{1}{\sqrt2}(a + b)
        \end{align*}

        $Prob(1) = |\braket{+}{\psi}|^2$ which is the squared magnitude of $\frac{1}{\sqrt2}(a+b)$. And the absolute square of a complex number $z$ is defined as $|z| = zz^*$ where $z^*$ is the complex conjugate.

        \begin{align*}
          Prob(1) &= \frac12 (|a|^2 + |b|^2 + ab^* + a^*b) \\
                  &= \frac12 (1 + ab* + a^*b) \\
                  &= \frac12 (1 + 2Re(ab^*)) \ a^*b \text{ is the complex conjugate of } ab^* \\
                  &= \frac12 + Re(ab^*)
        \end{align*}

        The second probability Prob($\frac12$) can be calculated like so:

        \begin{align*}
         \bra{-} &= \frac{1}{\sqrt2} (\bra{0} - \bra{1}) = \frac{1}{\sqrt2}(1 \ -1) \\
         \braket{-}{\psi} &= \frac{1}{\sqrt2}(1 \ -1) \begin{pmatrix} a \\ b \end{pmatrix} = \frac{1}{\sqrt2} (a -b)
        \end{align*}

        Using the same logic from Prob(1) where $\braket{-}{\psi}$ is just the squared magnitude of $\frac{1}{\sqrt2}(a-b)$, and the absolute square of a complex number $z$ is defined as $|z| = zz^*$ where $z^*$ is the complex conjugate, we can write out the following:

        \begin{align*}
          Prob(\frac12) &= \frac12 (|a|^2 + |b|^2 - ab^* - a^*b) \\
                  &= \frac12 (1 - (ab* + a^*b)) \\
                  &= \frac12 (1 - 2Re(ab^*)) \ a^*b \text{ is the complex conjugate of } ab^* \\
                  &= \frac12 - Re(ab^*)
        \end{align*}

        As a sanity check, if we add these two probabilities together, we should get 1:

        $$
        \frac12 + Re(ab^*) + \frac12 - Re(ab^*) = 1
        $$

      \end{enumerate}

    \item{Consider a state of two qubits $\ket{\psi} = \frac12 \ket{00} + \frac12 \ket{01} - \frac{i}{\sqrt2} \ket{11}$.}

      \begin{enumerate}
        \item{If you were to label the basis of 4-dimensional complex space in the following order: $\ket{00} \ket{01} \ket{10} \ket{11}$ what are the vector forms of $\ket{\psi}$ and $\bra{\psi}$?}

          For $\ket{\psi}$ the basis is defined by the coefficients which can be written in column vector form as such:

          \begin{equation*}
            \ket{\psi} = \begin{pmatrix} \frac12 \\ \frac12 \\ 0 \\ -\frac{i}{\sqrt2} \end{pmatrix} \\
          \end{equation*}

          For $\bra{\psi}$, this can be represented as the complex conjugate of the column vector rewritten as a row vector:

          \begin{equation*}
            \bra{\psi} = ( \frac12 \ \frac12 \ 0 \ \frac{i}{\sqrt2})
          \end{equation*}

        \item{If you were to label the basis of 4-dimensional complex space in the following order: $\ket{00} \ket{10} \ket{01} \ket{11}$ what are the vector forms of $\ket{\psi}$ and $\bra{\psi}$?}

          \begin{align*}
            \ket{\psi} &= \begin{pmatrix} \frac12 \\ 0 \\ \frac12 \\ -\frac{i}{\sqrt2} \end{pmatrix} \\
              \bra{\psi} &= ( \frac12 \ 0 \ \frac12 \ \frac{i}{\sqrt2})
          \end{align*}

        \item{If you were to label the basis of 4-dimensional complex space in the following order: $\ket{00} \ket{01} \ket{10} \ket{11}$ what is the matrix form of density operator $\rho = \ket{\psi} \bra{\psi}$?}

          $\rho$ is a $4 \times 4$ matrix formed by the outer product of $\ket{\psi}\bra{\psi}$.

          \begin{align*}
            \ket{\psi}\bra{\psi} &= \begin{pmatrix} \frac12 \\ \frac12 \\ 0 \\ -\frac{i}{\sqrt2} \end{pmatrix} \begin{pmatrix} \frac12 \ \frac12 \ 0 \ -\frac{i}{\sqrt2} \end{pmatrix} \\
                                 &= \begin{pmatrix}
                                   \frac14 & \frac14 & 0 & \frac{i}{2\sqrt2} \\
                                   \frac14 & \frac14 & 0 & \frac{i}{2\sqrt2} \\
                                   0 & 0 & 0 & 0 \\
                                   -\frac{i}{2\sqrt2} & -\frac{i}{2\sqrt2} & 0 & \frac12 \\
                                 \end{pmatrix}
         \end{align*}

        \item{Is $\ket{\psi}$ a separable state or entangled state?}

          $\ket{\psi}$ is separable if it can be written as $(a\ket{0} + b\ket{1}) \otimes (c\ket{0} + d\ket{1})$ for some $a,b,c,d$.

          For $\ket{\psi} = \frac12 \ket{00} + \frac12 \ket{01} - \frac{i}{\sqrt2} \ket{11}$, the coefficients are as follows:

          \begin{equation*}
            ac = \frac12, \ ad = \frac12, \ bc = 0, \ bd = -\frac{i}{\sqrt2}
          \end{equation*}

          With case $b = 0$, the equation $bd = -\frac{i}{\sqrt2} \ne 0$.

          With case $c = 0$, $ac = \frac12 \ne 0$.

          This state does not follow $(a\ket{0} + b\ket{1}) \otimes (c\ket{0} + d\ket{1})$, and therefore this state is entangled and cannot be separated.

      \end{enumerate}

    \item{Recall the quantum gate notation and diagrams, and consider the following circuit (not shown in my homework, but assume that it is there).}

      \begin{enumerate}
        \item{What is the output of the above circuit?}

          Applying Hadamard gate to the first qubit:

          \begin{equation*}
            \frac{1}{\sqrt2} (\ket{0} + \ket{1}) \otimes \frac{1}{\sqrt2} (\ket{0} + \ket{1})
          \end{equation*}

          Apply CNOT gate:

          \begin{align*}
            &= \frac{1}{2}[(\ket{0} + \ket{1}) \otimes (\ket{0} + \ket{1})] \\
            &= \frac{1}{2} (\ket{00} + \ket{01} + \ket{01} + \ket{11}) \\
            &= \frac{1}{2} (\ket{00} + \ket{01} + \ket{11} + \ket{10}) \ \text{CNOT applied}
          \end{align*}

          Apply phase change $e^{i\pi/4}$:

          \begin{align*}
            \frac{1}{2}(\ket{00} + \ket{10} + \ket{01} + e^{i\pi/4}\ket{11})
          \end{align*}

        \item{We measure the second qubit on the computational basis $\{\ket{0}, \ket{1}\}$ at the end of computation. After the measurement, what are the possible states of the first qubit, and the associated probabilities?}

          The equation is:

          \begin{equation*}
            \frac{1}{\sqrt2}(\ket{0}\ket{+} + \ket{1} \frac{1}{\sqrt2}(\ket{0} + e^{i\pi/4}\ket{1}))
          \end{equation*}

          For qubit 2 = 0:

          \begin{equation*}
            \frac12 \ket{00} + \frac12 \ket{10} = \frac12 (\ket{0} + \ket{1}) \otimes \ket{0} \\
          \end{equation*}

          For qubit 1 = 0:

          \begin{equation*}
            \frac12 \ket{01} + \frac12 e^{i\pi/4} \ket{11} = \frac12 (\ket{0} + e^{i\pi/4} \ket{1}) \otimes \ket{1}
          \end{equation*}

          The probability of an outcome = squared length ($\braket{\psi} = \sum |c_i^2|$)

          For qubit 2 = 0:

          \begin{equation*}
            |\frac12|^2 + |\frac12|^2 = \frac12
          \end{equation*}

          For qubit 2 = 1:

          \begin{align*}
            &|\frac12|^2 + |e^{i\pi/4}|^2 \\
            &=\frac14 + \sqrt{\cos(\pi/4)^2 + \sin(\pi/4)^2} \\
            &= \frac14 + \frac14 = \frac12
          \end{align*}

          Normalizing the branches:

          \begin{align*}
            k^2 \cdot \frac12 &= 1 \\
            k^2 &= 2 \\
            k &= \sqrt2 \\
            \frac{1}{\sqrt2}(\ket{0} &+ e^{i\pi/4}\ket{1}
          \end{align*}

          Apply to qubit 2 = 0:

          \begin{equation*}
            \sqrt2 (\frac12 \ket{0} + \frac12 \ket{1}) \\= \frac{1}{\sqrt2} (\ket{0} + \ket{1})
          \end{equation*}

          Apply to qubit 2 = 1:

          \begin{align*}
            \sqrt2 (\frac12 \ket{0} + \frac12 e^{i\pi/4} \ket{1}) \\= \frac{1}{\sqrt2} (\ket{0} + e^{i\pi/4} \ket{1})
          \end{align*}

          Measuring qubit 2 to get the outcome of 0: probability $\frac12$, and qubit 1 collapses to $\frac{1}{\sqrt2} (\ket{0} + \ket{1})$

          Measuring qubit 2 to get the outcome of 1: probability $\frac12$, and qubit 1 collapses to $\frac{1}{\sqrt2} (\ket{0} + e^{i\pi/4} \ket{1})$
          

      \end{enumerate}
    
\end{enumerate}

\end{document}
